% ch3 D.1 后段 (原书页 145–156 / PDF p160–p171)
% 衔接说明：本文件首段承上文件 ch3_D1b.tex 末句——原书 p144 末句 "The percentage
% is not in fact" 与本文件首词 “微不足道” 合成完整一句（“这个百分比实际上并不微不足道”）。
% 本文件止于原书 p156 末的 $\delta(q,q')$ 等式；原书 p157 顶部（D.2 小节标题之前）
% 尚有一段 D.1 收尾文字，已暂存于 ch3_D2.tex 顶部（%--D1-TAIL-- 标记之下），
% 装配时应接于本文件之后。

%JOIN%
微不足道，因为在较好的稀有气体和离子晶体中，$X^{2}T_{jk}$ 约为 1/2 至 1 eV，所以该振幅是 0.1 的量级。因此，令 $n_j = 0$ 确实是一个有实质内容的假定，不过在这个情形下也是一个相当准确的假定。

那么，就让我们在这个线性化近似下，把新的、完整的哈密顿量用玻色子算符 $\beta_{k\lambda}$ 表示出来（$\lambda$ 标记激子的偏振）：

\begin{gather*}
\mathcal{H} \;=\; \sum_{k\lambda} \left(E_0 + X^{2}\,T_{\lambda}(k)\right) \beta_{k\lambda}^{\dagger}\,\beta_{k\lambda} \\
+\;\sum_{k\lambda} X^{2}\,T_{\lambda}(k)\left(\beta_{k\lambda}^{\dagger}\,\beta_{-k\lambda}^{\dagger} + \beta_{k\lambda}\,\beta_{-k\lambda}\right) \tag{17}
\end{gather*}

让我们试着把它对角化。显然，唯一的途径是把算符 $\beta_k^{\dagger}$ 与 $\beta_{-k}$ 混合起来，反之亦然——这种做法读者如今已从超导理论中熟悉了。［其实，相应的变换很可能最早就是 Bogolyubov 在 $\mathrm{He}^{4}$ 理论中作出的 (87)。］于是我们令

\begin{gather*}
\beta_k^{\dagger} \;=\; u_k B_k^{\dagger} - v_k B_{-k} \qquad\qquad \beta_k \;=\; u_k B_k - v_k B_{-k}^{\dagger} \\
\beta_{-k}^{\dagger} \;=\; u_k B_{-k}^{\dagger} - v_k B_k \qquad\qquad \beta_{-k} \;=\; u_k B_{-k} - v_k B_k^{\dagger}
\end{gather*}

为保证这些新算符遵守玻色子的对易规则，我们只需检验

\begin{equation*}
\left[\,\beta_k,\;\beta_k^{\dagger}\right] = 1 \qquad\text{若}\qquad \left[\,B,\;B^{\dagger}\right] = 1
\end{equation*}

也就是说，

\begin{equation*}
u_k^{2}\left[\,B_k,\,B_k^{\dagger}\right] + v_k^{2}\left[\,B_{-k}^{\dagger},\,B_{-k}\right] = 1
\end{equation*}

或者，

\begin{equation*}
u_k^{2} - v_k^{2} = 1
\end{equation*}

只要令

\begin{equation*}
u_k \;=\; \cosh\frac{\theta_k}{2} \qquad\qquad v_k \;=\; \sinh\frac{\theta_k}{2}
\end{equation*}

它便自动得到满足。

现在让我们代入哈密顿量 (17) 之中。其中特定的一组项可以写成

\begin{equation*}
E_k^{0}\left(\beta_k^{\dagger}\beta_k + \beta_{-k}^{\dagger}\beta_{-k}\right) + E_k'\left(\beta_k^{\dagger}\beta_{-k}^{\dagger} + \beta_{-k}\beta_k\right)
\end{equation*}

这些项现在变为

\begin{align*}
E_k^{0}\,\Bigl\{\; & u_k^{2}\left(B_k^{\dagger}B_k + B_{-k}^{\dagger}B_{-k}\right) + v_k^{2}\left(B_{-k}B_{-k}^{\dagger} + B_k B_k^{\dagger}\right) \\
& -\; 2u_k v_k\,\left[B_k^{\dagger}B_{-k}^{\dagger} + B_{-k}B_k\right]\Bigr\} \\
+\; E_k'\,\Bigl\{\; & \left(u_k^{2} + v_k^{2}\right)\left(B_k^{\dagger}B_{-k}^{\dagger} + B_{-k}B_k\right) \\
& -\; u_k v_k\left(B_k^{\dagger}B_k + B_{-k}B_{-k}^{\dagger} + B_{-k}^{\dagger}B_{-k} + B_k B_k^{\dagger}\right)\Bigr\}
\end{align*}

这样，只要令

\begin{equation*}
\frac{2u_k v_k}{u_k^{2} + v_k^{2}} \;=\; \tanh\theta_k \;=\; \frac{E_k'}{E_k^{0}} \tag{18}
\end{equation*}

哈密顿量便被对角化了。

其余各项，在利用

\begin{equation*}
B_k B_k^{\dagger} \;=\; B_k^{\dagger} B_k + 1
\end{equation*}

之后，便可以写成

\begin{equation*}
\left[E_k^{0}\left(u_k^{2} + v_k^{2}\right) - 2u_k v_k E_k'\right]\left(n_k + n_{-k}\right) + 2\left(v_k^{2}E_k^{0} - u_k v_k E_k'\right)
\end{equation*}

于是，激子的新能量为

\begin{equation*}
E_k \;=\; E_k^{0}\cosh\theta_k - E_k'\sinh\theta_k
\end{equation*}

但根据 (18)，我们可以令

\begin{gather*}
E_k^{0} \;=\; E_k\cosh\theta_k \\
E_k' \;=\; E_k\sinh\theta_k
\end{gather*}

并可核实

\begin{equation*}
E_k \;=\; E_k\left(\cosh^{2}\theta_k - \sinh^{2}\theta_k\right)
\end{equation*}

进而得到

\begin{equation*}
E_k \;=\; \frac{E_k^{0}}{\cosh\theta_k} \;=\; E_k^{0}\left(1 - \tanh^{2}\theta_k\right)^{1/2} \;=\; \Bigl\{\left(E_k^{0}\right)^{2} - \left(E_k'\right)^{2}\Bigr\}^{1/2}
\end{equation*}

这一整套操作，与人们对角化 B.C.S. 哈密顿量 (88) 时所经历的步骤极为相似，只不过对玻色子来说，起作用的是双曲函数而不是三角函数，而且根号下的符号也相应地不同。Suhl（私人通信）曾作为一个有趣的练习，用玻色对算符（boson-pair operators）

\begin{equation*}
\beta_k^{\dagger}\beta_{-k}^{\dagger}\,, \qquad \beta_k^{\dagger}\beta_k + \beta_{-k}^{\dagger}\beta_{-k}\,, \qquad \beta_{-k}\beta_k
\end{equation*}

把这类哈密顿量对角化。这些玻色对算符具有双曲空间中而非实空间中的自旋所具有的性质——也就是说，它们像 $S$ 取虚值的自旋，对应于虚数阶的 Legendre 函数。

我们还看到，随着激子频率的降低，总能量也有所下降；这一下降可以被非常直接地解释为激子零点能相应的降低。利用恒等式

\begin{equation*}
v^{2} \;=\; \frac{u^{2} + v^{2} - 1}{2} \qquad\qquad \text{（由}\;\; u^{2} - v^{2} = 1\;\;\text{而来）}
\end{equation*}

我们得到

\begin{equation*}
\Delta E \;=\; \sum_{k}\left(E_k - E_k^{0}\right) \tag{19}
\end{equation*}

于是，经过完全对角化的总哈密顿量为

\begin{align*}
\mathcal{H} \;=\; \sum_{k>0,\,\lambda} \Bigl\{ & \left(E_0 + X^{2}T_{k\lambda}\right)^{2} - \left(X^{2}T_{k\lambda}\right)^{2}\Bigr\}^{1/2} \left(n_{k\lambda} + \frac{1}{2} + n_{k\lambda} + \frac{1}{2}\right) \\
& -\;\sum_{k\lambda}\left(E_0 + X^{2}T_{k\lambda}\right)
\end{align*}

（译注：第一个括号中第二个 $n_{k\lambda}$ 照录原书；按上下文应为 $n_{-k\lambda}$，如此对 $k>0$ 求和方能给出与 $k$ 与 $-k$ 两支模相应的零点项。）

其中 $n_{k\lambda} = B_{k\lambda}^{\dagger}B_{k\lambda}$。注意 $E_k = \left(E_0^{2} + 2E_0 X^{2}T_k\right)^{1/2}$。这显然与 (15$'$) 相同。

在这种表述中，我们使一个由 Hopfield 指出的相当有趣的事实变得一目了然：分子晶体的结合能可以重新解释为激子因偶极相互作用而导致的零点能之下降，因为它主要是 Van der Waals 吸引。

关于这个零点能问题，有一个相当重要的论点值得作些讨论。在我们原来的哈密顿量中，除了纯原子项 $(E_0/2)\left(c_p^{\dagger}c_p - c_s^{\dagger}c_s\right)$ 之外，其余各项都只涉及不同原子上的激子算符，而这些项当然是自动对易的。我们总是选择把它们写成 $b_j^{\dagger}b_k$ 的次序，结果得到的哈密顿量中含有 $E_k^{0}\,\beta_k^{\dagger}\beta_k$。这就是上面出现的那个减除项［即 (19) 中的 $-E_k^{0}$ 项］。当然，这一做法虽然是对我们实际物理情形的正确描述，却并不是经典介电连续介质在正确量子化下应有的哈密顿量；后者本应以对称化的形式出现：

\begin{equation*}
\frac{E_k^{0}}{2}\left(\beta_k^{\dagger}\beta_k + \beta_k\beta_k^{\dagger}\right)
\end{equation*}

由于我们所假设的系统与经典介电连续介质并没有明显的物理相似之处，这一点本来不必使我们过于担心；而事实上结果表明它无论如何也是对的。如果我们想采用对称化的形式，就必须补加 $\sum_k E_k^{0}$；鉴于 $\sum_k T_k = 0$，这个补项只依赖于纯原子的量 $E_0$，与相互作用毫无关系。当然，这个恒等式之所以成立，恰恰是因为我们本来就没有 $j=k$ 的相互作用项；这也正是我们一开始便能任意选择算符次序的原因。在自旋波问题和声子问题中，经典出发点与量子出发点之间的这种关系将更为切题。

现在让我们开始离开严格局域化的 Frenkel 激子这一主题，进而考虑这样的可能性：尽管受到对方 Coulomb 场的作用，受激电子或空穴偶尔仍会离开它的同伴而游走。这意味着我们必须在哈密顿量中计入如下各项：

\begin{equation*}
\mathcal{H}' \;=\; \sum_{jk}\left[b_{jk}^{s}\,c_j^{s\dagger}\,c_k^{s} + b_{jk}^{p}\,c_j^{p\dagger}\,c_k^{p}\right] \tag{20}
\end{equation*}

但同时我们还必须计入能量差 $U$；当受激电子与空穴处于同一原子上时，正是这个能量差使受激电子的能量降低：

\begin{equation*}
\mathcal{H}_{\mathrm{int}} \;=\; U\sum_j n_j^{p}\,n_j^{s} \tag{21}
\end{equation*}

对于 (20) 中的各项，我们所能采取的最好办法是用微扰论来处理：从 $U\to\infty$ 的极限出发，求出 $b^{2}/U$ 量级的第一级修正。在 U 非常大的极限下，合适的做法是考察 Coulomb 相互作用 (21) 的本征态。这些本征态按照同一原子上共有 0 对、1 对、2 对……s 电子与 p 电子而自行形成一个分层结构，其能量分别为 0、U、2U……迄今为止，我们处理的只是这类态中 $U = 0$ 的简并子空间，即电子与空穴占据同一原子的那些态；这些态已被偶极相互作用劈开，而这一劈裂与 U 相比显然很小。下一步要计算的，是这个简并子空间因虚跃迁而产生的进一步劈裂；在虚跃迁中，电子或空穴会跳到一个已经有另一粒子所在的原子上，这在 U 非常大的极限下具有能量 U。显然，为了得到最低阶的效应，我们不能容许电子继续游荡，而必须强迫它在下一次跳跃时就回到自己的空穴处。于是，我们把微扰级数

\begin{equation*}
\mathcal{H}_{\mathrm{eff}} \;=\; \mathcal{H}' - \mathcal{H}'\,\frac{1}{U}\,\mathcal{H}' + \mathcal{H}'\,\frac{1}{U}\,\mathcal{H}'\,\frac{1}{U}\,\mathcal{H}' \cdots
\end{equation*}

投影到由“允许”态（即电子与空穴在一起的态）构成的子空间 $|0\rangle$ 上：

\begin{equation*}
\mathcal{H}^{(2)} \;=\; -\left\langle 0\right|\,\mathcal{H}'\,\frac{1}{U}\,\mathcal{H}'\,\left|0\right\rangle
\end{equation*}

并略去一切更高阶的微扰。只要只保留电子或空穴返回其原来所在原子的那些项，这一点便很容易做到；(20) 与 (21) 的二级效应于是为

\begin{align*}
\mathcal{H}^{(2)} \;=\; -\frac{1}{U}\sum_{jk}\Bigl[\; & b_{jk}^{s}\left(c_j^{s}\right)^{\dagger}c_k^{s} + b_{jk}^{p}\left(c_j^{p}\right)^{\dagger}c_k^{p}\,\Bigr] \\
\times\,\Bigl[\; & b_{kj}^{s}\left(c_k^{s}\right)^{\dagger}c_j^{s} + b_{kj}^{p}\left(c_k^{p}\right)^{\dagger}c_j^{p}\,\Bigr] \\
=\; -\frac{1}{U}\sum_{jk}\Bigl\{\; & |b_{jk}^{s}|^{2}\,n_j^{s}\left(1 - n_k^{s}\right) + |b_{jk}^{p}|^{2}\,n_j^{p}\left(1 - n_k^{p}\right) \\
& +\; b_{jk}^{s}\,b_{kj}^{p}\left(c_j^{s}\right)^{\dagger}c_k^{s}\left(c_k^{p}\right)^{\dagger}c_j^{p} \;+\; \mathrm{c.c.}\Bigr\}
\end{align*}

（译注：第二个方括号内 $\left(c_k^{p}\right)$ 上的产生算符符号 $^{\dagger}$ 原书漏印，此处按对称性补出。）

前两项有些不对称，所以我们在 jk 与 kj 这两对指标之间把它们对称化。同时我们还考虑到如下事实：在我们所考虑的 $|0\rangle$ 态流形之内，$n_j^{p} = 1 - n_j^{s}$，反之亦然——每个原子上要么有一个受激电子，要么有一个基态电子。于是，对这两项我们得到

\begin{equation*}
-\frac{1}{U}\sum_{jk}\frac{|b_{jk}^{s}|^{2} + |b_{jk}^{p}|^{2}}{2}\left(n_j^{s}n_k^{p} + n_j^{p}n_k^{s}\right)
\end{equation*}

用激子算符来表示，

\begin{align*}
n_j^{s}n_k^{p} + n_j^{p}n_k^{s} \;&=\; \frac{1}{2}\left(n_j^{s} + n_j^{p}\right)\left(n_k^{s} + n_k^{p}\right) - \left(n_j^{p} - n_j^{s}\right)\left(n_k^{p} - n_k^{s}\right) \\
&=\; \frac{1}{2}\left(1 - \sigma_z^{j}\sigma_z^{k}\right)
\end{align*}

或者

\begin{equation*}
=\;2\left(-\,n_j n_k + n_j + n_k\right)
\end{equation*}

后一种形式表明，这一项除了给出相邻激子之间的一个非线性相互作用之外，还使激子的能量降低

\begin{equation*}
\frac{1}{U}\sum_{K}\left\{|b_{jk}^{s}|^{2} + |b_{jk}^{p}|^{2}\right\}
\end{equation*}

（译注：求和限 $K$ 原书如此（大写），按上下文应为 $jk$。）

$\sigma$ 形式在自旋波问题中将会很重要。

激子能量中的这个对角项并不十分重要；它所反映的事实是：$b$ 项把自由粒子的各个能级展宽成为一个能带，从而把激子的平均能量压低。另一项则更有意思；它恰好具有导致激子运动的那些偶极项的形式，因此我们可以把这些项写成与偶极项相类似的形式：

\begin{gather*}
\mathcal{H}^{(2)} \;\simeq\; -\sum_{jk} T_{jk}'\,b_j^{\dagger}\,b_k \\
T_{jk}' \;=\; \frac{b_{jk}^{s}\,b_{jk}^{p}}{U} \tag{22}
\end{gather*}

这里涉及的机制显然是：电子先虚跳到相邻的原子上，随后空穴再朝同一方向跳过去。这种迁移机制对光学活性的 Frenkel 激子并不十分重要，但对禁戒的激子却相当重要，对自旋波尤其如此。显然，这一效应的效果是使激子谱略微跟上由 $b_{jk}$ 引起的单粒子谱的展宽（图 46）。一旦 b 大到足以与 U 相比，这个模型就完全失效了。当 b 增大时我们看到，激子态的能量减小约 $Zb^{2}/U$；另一方面，最低的一些能带态——也就是最能利用与近邻之间正确相位关系的那一些态——其能量则减小约 $Zb$。于是，在这个模型中激子的结合能按如下方式减小：

\begin{equation*}
\mathrm{BE} \;\cong\; U - Zb + \frac{Zb^{2}}{U}
\end{equation*}

%FIG{46}: 图 46　能级示意：p 能级与激子（Exciton）能级各给出 no b（不计 $b$，水平直线）与 with b（计入 $b$，曲线）两种情形；$U$ 标出 p 组态与激子组态之间的能量间隔，$E_0$ 为自 s 能级到激子能级的间距，左侧箭头 $E$ 指示能量向下减小，横轴为 $k$（原书图注仅作 “Figure 46”）

因此，一旦 $bZ \sim U$，纯局域激子的能量就未必低于带边（图 47）；那时我们将不得不只用能量极小点附近的电子态来构成激子态，才有希望得到低于能带连续谱的激子。此外，在一个 b 相当可观的材料中，Wannier 函数更为扩展，而介电屏蔽会使 U 的作用减弱。于是我们预期，随着 $b/U$ 的变化，会发生一次相当陡峭的转变，转到激子理论的相反情形，即 Wannier 极限：在这种极限下，电子与空穴很少同时处于同一原子上，而是具有一个可以由能带极小附近的态构成的长程波函数。

%FIG{47}: 图 47　向下弯曲的 “Band”（能带）曲线与 “Localized exciton”（局域激子）水平能级线相交示意（原书图注仅作 “Figure 47”）

这种情形可以用能带理论的“有效单带哈密顿量”（effective one-band hamiltonian）方法来处理——事实上，Wannier 在 1937 年正是从这个问题出发开创了这一理论的近代发展 (78)。我们采用近似哈密顿量

\begin{equation*}
\mathcal{H}_{\mathrm{eff}} \;=\; E_{\mathrm{el}}(k_e) + E_{\mathrm{hole}}(k_h) + \frac{e^{2}}{\kappa\,|R_e - R_h|}
\end{equation*}

一般说来，这会导致相当复杂的方程，尽管原则上并非不可处理。在简并能带的情形下尤其如此，那时甚至连质心运动也不容易分离出来。幸运的是，最有趣的那些情形恰好也是激子相当容易被观察到的情形，而这些情形看来都是简单的有效质量椭球（effective mass ellipsoid）情形；在这种情形下我们可以写出

\begin{equation*}
\mathcal{H} \;=\; -\frac{\hbar^{2}}{2}\sum_{\alpha=x,y,z}\frac{1}{m_{\alpha e}}\frac{\partial^{2}}{\partial x_{\alpha e}^{2}} + \frac{1}{m_{\alpha h}}\frac{\partial^{2}}{\partial x_{\alpha h}^{2}} - \frac{e^{2}}{\kappa\,|R_e - R_h|}
\end{equation*}

并且质心运动与相对运动可以分离开来，给出描写整体运动的激子质量张量 $\left(m_e + m_h\right)_{\alpha}$，以及描写相对运动的约化质量张量

\begin{equation*}
\left(\frac{m_e m_h}{m_e + m_h}\right)_{\alpha}
\end{equation*}

即使这一点做不到——例如出现简并，或者椭球的主轴互不相同——往往也会有一个粒子比另一个重得多，从而可以找到良好的近似。

于是，激子在坐标表象中的波函数为

\begin{equation*}
\Psi \;=\; \Phi(R_e - R_h)\;e^{iq\left(R_e + R_h\right)}
\end{equation*}

其中 q 是质心运动的波数。相对于原子间距而言，$\Phi$ 是一个相当长程的函数。把这一处理中所作的近似与我们前面 Frenkel 处理中的近似作一对比，并用图示来说明激子在这两种情形下如何运动，是有趣的。让我们把两个能带——电子（导）带与空穴（价）带——中 Wannier 函数的产生算符分别记为 $w_e^{\dagger}(R_j)$ 与 $w_h^{\dagger}(R_j)$；例如，$w_e$ 定义为

\begin{equation*}
w_e^{\dagger}(\mathbf{R}_j) \;=\; \int a_e(r - \mathbf{R}_j)\,\psi^{\dagger}(r)\;dr
\end{equation*}

这里 e 表示空带。波函数 $\Psi$ 可以写成

\begin{equation*}
\Psi \;=\; \sum_{jk}\phi_q(R_j - R_k)\;e^{iq\cdot\left(R_j + R_k\right)}\;w_e^{\dagger}(\mathbf{R}_j)\,w_h(\mathbf{R}_k)\,\Psi_g
\end{equation*}

而这个近似哈密顿量相当于在相互作用能中只保留

\begin{equation*}
\frac{e^{2}}{\kappa\,|R_j - R_k|}\;n_e(R_j)\,n_h(R_k)
\end{equation*}

并忽略一切非对角元。于是在这种情形下，相互作用根本不引起任何运动：$w_e^{\dagger}(j)\,w_h(k)\,\Psi_g$ 就是这一相互作用的一个本征函数，其能量为

\begin{equation*}
-\frac{e^{2}}{\kappa\,|R_j - R_k|}
\end{equation*}

运动反而完全来自哈密顿量中的 $b_{jk}$ 各项。我们可以用一幅以时间 t 对位置 x 作图的图来描述它：相互作用发生在电子与空穴位置固定之处，而两次相互作用之间这两个粒子在 x 空间中运动（图 48）。这些相互作用只在长距离处发生，其唯一的效果就是把电子与空穴维系在一起。

%FIG{48}: 图 48　时间 $t$–位置 $x$ 图：电子（$k_e$）与空穴（$k_h$）的径迹在固定位置 $R_j$、$R_k$ 之间由一条波线相连，波线标 “Interaction”（相互作用）$-e^{2}/\kappa|R_j - R_k|$，作用后两径迹以 $k_e'$、$k_h'$ 继续（原书图注仅作 “Figure 48”）

在另一种情形中，我们假定 $b_{jk}$ 为零，电子与空穴完全不能移动；唯一的例外是 $w_h$ 与 $w_e^{*}$ 可以彼此湮灭，其矩阵元为 X［见图 49(a)］。图 49(b) 中画出的，是由式 (22) 所描述的经由电子或空穴一次虚跳的运动；它与图 48 的过程密切相关。图 49(a) 所示的那种图也会在 Wannier 激子中出现，但只在电子与空穴的路径偶然相交这一罕见事件中才会发生（图 50）。该图表明这样一个事实：与外场相互作用的有效偶极矩——其图解如图 51 所示——要被缩小为激子态体积的倒数，因为只有当电子与空穴彼此靠近时湮灭才会发生。由此造成的纵–横频率差相当小。它无法单靠有效哈密顿量理论来处理。

%FIG{49}: 图 49　$t$–$x$ 图两种情形：(a) 电子与空穴的径迹在 $R_j$ 与 $R_k$ 处各标一次 “X”，其间由标 “T” 的波线相连，两端各有向上的箭头；(b) 虚跳图：标 $b_{jk}$ 与 $b_{kj}$ 的竖直跳变跨过标 $U$ 的水平线（原书图注仅作 “Figure 49”）

%FIG{50}: 图 50　Wannier 激子中电子与空穴路径偶然交叉时的相互作用图（$t$–$x$ 图，交叉点处接波线）（原书图注仅作 “Figure 50”）

%FIG{51}: 图 51　与外场相互作用的顶点图：一条波线（外场量子）在同一个时空点上与电子、空穴的径迹相接（原书图注仅作 “Figure 51”）

看待我们对基态能量以及 Frenkel 激子频率所作的那些零点修正的一种方式，是注意到它们相当于在激子的传播方程中计入了所谓的“反向图”（backwards diagrams）（见图 52）。这类“反向图”在核多体理论和核能级结构中的效应要大得多 (89)。这些图例示了在多体理论中对图的子集求和的过程。

%FIG{52}: 图 52　“反向图”示意：电子与空穴的径迹在 $R_i$、$R_j$、$R_k$、$R_l$ 各处经波线相互作用，中间形成两个环状结构，各段标以 $b_i$、$b_j$、$b_j^{\dagger}$、$b_k$、$b_k^{\dagger}$、$b_l^{\dagger}$（原书图注仅作 “Figure 52”）

像图 48 那样的图称为“梯图”（ladder diagrams）；对两个粒子求解 Schrödinger 方程，就等价于“把梯图求和到无穷阶”——也就是说，允许这种隔距相互作用按需要的次数、实质上连续不断地起作用。像图 49 那样的图则常称为“泡图”（bubble）或“环图”（ring），它们与材料的介电常数有关。在激子这个简单情形中，除了某种程度的数学复杂性之外，其实没有什么东西妨碍我们把两组图一起求和；我们已经给出了两类求和分别更为重要的那些极限情形。在更复杂的集体激发情形中，也许只能单独地“求和”梯图或泡图。在“泡图”的情形下，通常还必须把“反向”图也一并求和；这一步由无规相位近似（random-phase approximation）完成 (90)。

对 Wannier 型的激子来说，没有什么特别明显或简单的方式可以直接过渡到玻色子算符。相反，我们必须遵循较为就事论事的（ad hoc）做法：写下激子算符，然后直接验证，在\emph{只有很少数激发被占据的态}中，它们的对易规则近似地就是玻色子的对易规则。产生一个 Wannier 激子的算符可以写成

\begin{equation*}
O^{\dagger}(q) \;=\; \sum_{jk}\phi_q(j-k)\;e^{iq\cdot\left(\mathbf{R}_j + \mathbf{R}_k\right)}\;w_e^{\dagger}(\mathbf{R}_j)\,w_h(\mathbf{R}_k)
\end{equation*}

于是我们得到

\begin{align*}
\left[O(q),\,O^{\dagger}(q')\right] =\; & \sum_{jklm}\phi_q(j-k)\,\phi_{q'}(l-m)\;e^{iq\cdot\left(\mathbf{R}_j - \mathbf{R}_k\right)\,-\,iq'\cdot\left(\mathbf{R}_l + \mathbf{R}_m\right)} \\
\times\; & \left[\,w_h^{\dagger}(\mathbf{R}_l)\,w_e(\mathbf{R}_m),\;w_e^{\dagger}(\mathbf{R}_j)\,w_h(\mathbf{R}_k)\,\right]
\end{align*}

把这个对易子作用到基态 $\Psi_g$ 上时，除非 $R_j = R_m$、$R_k = R_l$，否则它将为零，因为基态中\emph{从一开始}就不存在任何激子。这意味着 $n_h = 1$、$n_e = 0$；也就是说，除非 l 上有一个空穴、\emph{且} m 上有一个受激电子，否则 $w_h^{\dagger}\,w_e = 0$。在该情形下，对易子的值为 1。在激子情形中，甚至从物理上讲，人们也很少会达到使这一假设失效的态，或许在某些光学微波激射器（optical maser）中除外。于是

\begin{align*}
\left[O(q),\,O^{\dagger}(q')\right]\;&=\;\sum_{j,k}\phi_q(R_j - R_k)\,\phi_{q'}(R_j - R_k)\;e^{i(q-q')\cdot\left(R_j + R_k\right)} \\
&=\;\delta(q,\,q')
\end{align*}
